% Section 8 (part 6 of 8): the model identity; signs, weights and identification; the relabelling.

\begin{lem}[The model identity]
\label{lem:con-S2}
Let $W_1(\tilde p):=\tilde p^4/4+1/\tilde p$. For every $b\in\bbZ/5$, $\SPn_{\tilde p=1}[-W_1;\tilde p^{\nu_b}](t)=\varepsilon_b\Fs_b(t)$; equivalently, the full stationary phase $\mathrm{SP}_{\tilde p=1}[\tilde p^{\nu_b}\re^{-W_1(\tilde p)/t}\,\dd\tilde p]$ divided by $(2\pi t/5)^{1/2}\re^{-5/(4t)}$ is $\varepsilon_b\Fs_b(t)$.
\end{lem}

\begin{proof}
(The letters $J_m$, $G_b$, $U$ are local.) Let $A:=\bbQ\fps{v-1}$ and $\varsigma:=v^{1/5}\in A$. The phase $-\Wv(p)$ has the critical point $p=\varsigma$ (since $\Wv'(p)=p^3-v/p^2$) with $\Wv''(\varsigma)=5\varsigma^2$, a unit of $A$. Let $J_m:=\mathrm{SP}_{p=\varsigma}[p^m\re^{-\Wv(p)/z}\,\dd p]$, the full stationary phase over $A$. By \cref{lem:fsp-dq}, with the moving critical point $\varsigma$, $z\,\dd J_m/\dd v=-J_{m-1}$; by \cref{lem:fsp-ibp} applied to $\dd(\re^{-\Wv/z})$, $J_3=vJ_{-2}$. Hence $(J_0,-J_{-1},-J_3,J_2,J_1)=(\varepsilon_bJ_{\nu_b})_b$ solves the model system $z\,\dd Y_b/\dd v=b^{\mathrm m}_bY_{b+1}$, $b^{\mathrm m}=(1,-1/v,1,-1,-1)$ (\cref{prop:set-modrows}); the row $b=1$ uses $J_{-2}=J_3/v$, the other rows the derivative identity alone. Next, \cref{lem:fsp-cov} over $A$ with the linear map $p=\varsigma\tilde p$, for which $\Wv(\varsigma\tilde p)=\varsigma^4W_1(\tilde p)$, followed by \cref{lem:fsp-homog} with the unit $\mu=\varsigma^4$, gives
\[
J_{\nu_b}=(2\pi z/5)^{1/2}\re^{U/z}\varsigma^{\nu_b-1}G_b(z/\varsigma^4),\qquad U=-\tfrac54\varsigma^4,
\]
where $G_b(t):=\SPn_{\tilde p=1}[-W_1;\tilde p^{\nu_b}](t)$ does not depend on $v$. With $\dd U/\dd v=-1/\varsigma$ and $\dd\varsigma/\dd v=1/(5\varsigma^4)$ the chain rule gives $z\frac{\dd}{\dd v}\bigl[\re^{U/z}\varsigma^{k}F(z/\varsigma^4)\bigr]=\re^{U/z}\varsigma^{k-1}\mathcal T_kF$ with $\mathcal T_k=-1+\frac t5(k-4t\frac{\dd}{\dd t})$, so the system becomes the colour-$0$ T-system of \cref{prop:set-modrows} for $F_b:=\varepsilon_bG_b$. Since $G_0(0)=1$, the uniqueness half of \cref{prop:set-modrows} (pivot $4n$) gives $F_b=\Fs_b$, i.e. $G_b=\varepsilon_b\Fs_b$.
\end{proof}

\begin{lem}[Signs, weights, identification]
\label{lem:con-id}
Work in a conifold frame of \cref{def:set-frames}, and let $Y_c=\tau^{-1/4}$ and the $\pi_m$ be as in \cref{lem:con-cfu,lem:con-red}.
\begin{enumerate}
\item[(a)] $\mathrm{SP}_{Y_c}[Y^{\nu_b}\re^{-W_V/\zh}\,\dd Y]=N_{\zh}\,\varepsilon_b\,\tau^{(1-\nu_b)/4}\Fs_b(\tau\zh)$, where $N_{\zh}:=(2\pi\zh/5)^{1/2}\re^{-W_V(Y_c)/\zh}$ and $W_V(Y_c)=\hat\cW(1)$.
\item[(b)] $\PhiSP=\sum_b\varepsilon_b\,\pi_{\nu_b}\,\tau^{(1-\nu_b)/4}\Fs_b(\tau\zh)$.
\item[(c)] Let $P$ be a graded regular gauge from the pulled-back model to the quintic system (\cref{def:gau-onepole,prop:gau-graded}), and let $\Phi_P$ be $\varpi_0$ times the colour-$0$ entry of the hatted row $0$ of $P\,\Qr^{\mathrm m}(v(u))$; for $P=\Pc$ this is $\Phican$ (\cref{thm:gau-R0}). Then
\[
\Phi_P/\varpi_0=\sum_b\tau^{(1-\nu_b)/4}\,G_b\,\Fs_b(\tau\zh),\qquad G_b:=s^{\nu_b}P_{0b}(u,s^4\zh),
\]
where $G_b\in\bbQ\fps{u}\fps{\zh}$ contains only powers $\zh^k$ with $k\equiv-b\bmod5$, and $[\zh^k]G_b$ is divisible by $u^{(\nu_b+4k)/5}$.
\item[(d)] Let $\Delta_b\in\bbQ\fps{u}\fps{\zh}$, $b\in\bbZ/5$, be series that contain only powers $\zh^k$ with $k\equiv-b\bmod5$. If $\sum_b\tau^{(1-\nu_b)/4}\Delta_b\Fs_b(\tau\zh)=0$, then every $\Delta_b$ vanishes.
\end{enumerate}
Hence, if $P$ is a graded regular gauge with $\Phi_P=\PhiSP$, then $G_b=\varepsilon_b\pi_{\nu_b}/\varpi_0$ for every $b$.
\end{lem}

\begin{proof}
(a) Put $Y=Y_c\tilde p$. Since $V=Y_c^5$, $W_V(Y_c\tilde p)=Y_c^4W_1(\tilde p)$, and $\zh/Y_c^4=\tau\zh$. By \cref{lem:fsp-cov} (a linear map over $\bbQ\fps{u}$) and \cref{lem:fsp-homog} with $\mu=Y_c^4=\tau^{-1}$, the left side is $Y_c^{\nu_b+1}\,\mathrm{SP}_{\tilde p=1}[\tilde p^{\nu_b}\re^{-W_1/t}\,\dd\tilde p]$ at $t=\tau\zh$, which by \cref{lem:con-S2} is $Y_c^{\nu_b+1}(2\pi\tau\zh/5)^{1/2}\re^{-5/(4\tau\zh)}\varepsilon_b\Fs_b(\tau\zh)$. Here $Y_c^{\nu_b+1}\tau^{1/2}=\tau^{(1-\nu_b)/4}$ and $\frac5{4\tau}=\frac54Y_c^4=W_V(Y_c)=\hat\cW(1)$ (\cref{lem:con-cfu}).

(b) By \cref{lem:fsp-cfu} the full stationary phase at $X=1$ of $\re^{-\hat\cW/\zh}\hat\fa\,\dd X$ equals that at $Y_c$ of its image $\re^{-W_V/\zh}\fa_Y\,\dd Y$. By \cref{lem:con-red} and \cref{lem:fsp-ibp}, applied once to the exact form $\dd(\mathcal E\re^{-W_V/\zh})$, whose primitive $\mathcal E$ is one amplitude over $\bbQ\fps{u}$, and by the $\bbQ\fps{u}\fps{\zh}$-linearity of the functional (\cref{def:fsp-wick}), this is $\sum_m\pi_m\,\mathrm{SP}_{Y_c}[Y^m\re^{-W_V/\zh}\,\dd Y]$. The exponents $m=-1,\dots,3$ are the $\nu_b$; insert (a) and divide by $N_{\zh}=(2\pi\zh/5)^{1/2}\re^{-\hat\cW(1)/\zh}$ (\cref{def:con-PhiSP}).

(c) The colour-$0$ model column is $\Qr^{\mathrm m}_b=\varsigma^{\nu_b-1}\Fs_b(z/\varsigma^4)$ after the common exponential is stripped (\cref{prop:set-modrows,thm:gau-R0}), and $\varsigma=s\tau^{-1/4}$, $z/\varsigma^4=\tau\zh$ (\cref{def:set-tau}). Row $0$ of $P\Qr^{\mathrm m}$, hatted by $s^{1-\nu_0}=s$, is $\sum_bs\,P_{0b}\,s^{\nu_b-1}\tau^{(1-\nu_b)/4}\Fs_b(\tau\zh)$, which is the stated formula. For graded $P=\sum_kz^k\sum_{b'}p_{k,b'}E_{b'+k,b'}$ the entry $P_{0b}$ contains only $z^k$ with $k\equiv-b$, and $s^{\nu_b}z^k=u^{(\nu_b+4k)/5}\zh^k$ with $\nu_b+4k\equiv-b-k\equiv0\bmod5$ and $\nu_b+4k\geq0$.

(d) If not all $\Delta_b$ vanish, let $N$ be the least $\zh$-degree occurring in them. Only $b\equiv-N$ can contribute to $\zh^N$, through $\tau^{(1-\nu_b)/4}\varepsilon_b[\zh^N]\Delta_b$, because $\Fs_b(0)=\varepsilon_b$; as $\tau$ is a unit and $\varepsilon_b=\pm1$, this coefficient vanishes, a contradiction.

Finally, if $\Phi_P=\PhiSP$, subtract (b) from $\varpi_0$ times (c): the family $\Delta_b:=G_b-\varepsilon_b\pi_{\nu_b}/\varpi_0$ is graded as in (d), by (c) and \cref{lem:con-red} ($\nu_b\equiv-b$), so it vanishes.
\end{proof}

\begin{lem}[The relabelling]
\label{lem:con-relab}
Let $\hat\cV:=\Res_X[\re^{-\hat\cW/\zh}\hat\fa\,\dd X]$, computed coefficientwise in $u$ in the residue ring $\cR_{\mathrm{res}}$ of \cref{lem:fsp-res}; at each $u$-order the integrand is $\re^{-(1/X+X^4/4)/\zh}$ times a Laurent polynomial in $X$ over $\bbQ[\zh,\zh^{-1}]$, an element of $\cR_{\mathrm{res}}$. Then $\cV=s\,\hat\cV$ coefficientwise: $[u^N\zh^j]\hat\cV=[u^{N+(1-4j)/5}z^j]\cV$ for all $N$ and $j$.
\end{lem}

\begin{proof}
\emph{Labels.} Monomials are $u^Ez^\beta x^\alpha$ in the flat variables and $u^N\zh^jX^\alpha$ in the hatted ones; the substitution $x=sX$, $z=s^4\zh$, $u=s^5$ acts on exponents by $\iota(E,\beta,\alpha)=(E+(4\beta+\alpha)/5,\beta,\alpha)$, which is injective. With $\dd x=s\,\dd X$ the claim is the coefficientwise statement.
\emph{The family.} The logarithm of $\re^{-\cG/z}\fa$ is a sum of atoms $-x^{5m-1}/(m(5m-1)z)$, $x^{5n}/(2n)$ and $\bk_k\,c_{k,P}\,(xz)^{2k-1}x^{5P}$, with $c_{k,P}:=5^{1-2k}\delta_{P,0}-5\binom{2k-2+5P}{5P}$ (\cref{lem:con-form}~(b)). A member of the family is a triple ($k$, $r$, a finite multiset of atoms); its term is $[u^r](\Xi/u)^k\,\frac{(-1)^k}{k!}\,u^{k+r}(xz)^{-k}$ times the product of the atoms divided by the factorials of the multiplicities. By \cref{def:con-PhiSP}, $\cV$ is the sum of the $x^{-1}$-parts of these terms.
\emph{Flat fibres.} A fixed $(u^E,z^\beta,x^{-1})$ receives finitely many members: $k\leq E$, and every atom has positive $x$-degree, so the multiset is bounded by the $x$-degree $k-1$.
\emph{Hatted fibres.} Every atom $x^az^b$ has hatted $u$-order $(a+4b)/5\geq0$, the only atom of hatted order $0$ is $-x^4/(4z)\mapsto-X^4/(4\zh)$, and $u^{k+r}(xz)^{-k}\mapsto u^r(X\zh)^{-k}$. Fix $(u^N,\zh^j,X^{-1})$. Then $r\leq N$, at most $N$ atoms have positive hatted order, from a finite list, and the $X$- and $\zh$-balances $-k+4l+a'=-1$, $-k-l+b'=j$, where $l$ is the multiplicity of $X^4/(4\zh)$ and $a'$, $b'$ are the degrees of the other atoms, determine $k$ and $l$ (the determinant is $5$). So this fibre is finite too.
\emph{$\hat\cV$ is the hatted sum.} In $\cR_{\mathrm{res}}\fps{u}$ the product is the convergent sum of the products of coefficients, and $\re^{-(\Xi/u)/(X\zh)}=\re^{-1/(X\zh)}\re^{-(\Xi/u-1)/(X\zh)}$ holds coefficient by coefficient: at $(u^N,X^{-n})$ both sides equal $(-1/\zh)^n[u^N](\Xi/u)^n/n!$, a finite expression. So $\hat\cV$ is the sum of the family in the hatted order.
\emph{Conclusion.} $\cV$ and $s\hat\cV$ are sums of one family over the fibres of two injective labellings related by $\iota$, with finite fibres in both orders; so they agree coefficient by coefficient.
\end{proof}

% Section 8 (part 7 of 8): the key lemma, the weight lemma, the residue, Theorem P(b).

\begin{lem}[The key lemma]
\label{lem:con-key}
Expand $\Sigma:=\sum_b\pi_{\nu_b}\,s^{-\nu_b}\,\IC_{\nu_b}(v(u),z)$ into monomials $u^Ez^\beta$, writing $\zh=s^{-4}z$ in $\pi_m$,
\[
\IC_m(v,z)=\sum_{i'\geq0}\frac{(-1/(4z))^{i'}}{i'!}\,\frac{(-v/z)^e}{e!},\qquad e:=m+4i'+1\quad(\text{terms with }e\geq0),
\]
and $v^e=u^eV^e$ (\cref{def:gau-van,def:set-tau}). A monomial $u^i\zh^j$ of $\pi_m$, the $i'$-th term of $\IC_m$ and a monomial $u^l$ of $V^e$ contribute to $u^Ez^\beta$ with $\beta=j-5i'-m-1$ and $E=i+l-(4\beta-1)/5$. Hence only $\beta\equiv4\bmod5$ occurs, and the coefficient of $u^Ez^\beta$ receives contributions only from the $u^i$-coefficients of the $\pi_m$ with $i\leq E+(4\beta-1)/5$, through finitely many terms; $\Sigma$ is a well-defined formal sum of monomials $u^Ez^\beta$.
\end{lem}

\begin{proof}
The $z$-exponent is $j-i'-e=j-5i'-m-1$, and $j\equiv m\bmod5$ by \cref{lem:con-red}, so $\beta\equiv-1$. The $u$-exponent is $i-\frac{4j}5-\frac m5+e+l$; substituting $j=\beta+5i'+m+1$ and $e=m+4i'+1$ gives $E=i+l-(4\beta-1)/5$. So $i=E+(4\beta-1)/5-l\leq E+(4\beta-1)/5$; for fixed $(E,\beta)$ and $i$ the exponent $l$ is fixed, $\pi_m$ is a polynomial in $\zh$ at the order $u^i$, so finitely many $j$ occur, and then $i'$ is fixed.
\end{proof}

Finite fibres give a formal sum. That $\Sigma$ lies in $\bbQ\lser{z}\fps{u}$ follows from \cref{prop:con-res} below, which identifies it coefficientwise with $\cV\in z^{-1}\bbQ[z^{-1}]\fps{u}$, or directly from the following lemma.

\begin{lem}[The weight lemma]
\label{lem:con-weight}
With the weight $\omega(u^NY^n\zh^j):=5N-n-4j$, every monomial $u^i\zh^jY^m$ of $\pi_mY^m$ has $\omega\geq0$; hence $E\geq0$ in \cref{lem:con-key}, and $\Sigma\in\bbQ\lser{z}\fps{u}$.
\end{lem}

\begin{proof}
A monomial $x^az^b$ of $\fa$ becomes $u^{(a+4b)/5}X^a\zh^b$, of weight $0$. The uniformising map and its inverse are $X(1+\sum c\,u^kX^{5d})$ and $Y(1+\sum c'\,u^kY^{5d})$ with $d\leq k$ (\cref{lem:con-cfu}), so they do not lower the weight, and every monomial of $\fa_Y$ has weight at least $0$. The step $Y^n\mapsto(n-3)\zh Y^{n-4}$ of the reduction keeps the weight, the step $Y^n\mapsto V_lu^lY^{n-5}$ raises it by $5(l+1)$. Hence $5i-m-4j\geq0$, and $E=(5i-4j-m)/5+e+l\geq0$; for fixed $E$, $\beta\geq(1-5E)/4$.
\end{proof}

\begin{prop}[The residue]
\label{prop:con-res}
Work in a conifold frame, and let $P$ be a graded regular gauge with $\Phi_P=\PhiSP$ (\cref{lem:con-id}). Then
\[
\cV=\varpi_0\,\bigl(P\,\YC(v(u),z)\bigr)_0\qquad\text{in }\bbQ\lser{z}\fps{u}.
\]
\end{prop}

\begin{proof}
By \cref{lem:con-relab}, $\cV=s\hat\cV$ coefficientwise. By \cref{lem:fsp-res} the residue on $\cR_{\mathrm{res}}\fps{u}$ is invariant under the uniformising map, termwise in $u$, and vanishes on exact forms; with \cref{lem:con-red},
\[
\hat\cV=\Res_Y\bigl[\re^{-W_V/\zh}\fa_Y\,\dd Y\bigr]=\sum_{m=-1}^3\pi_m\,\Res_Y\bigl[Y^m\re^{-W_V/\zh}\,\dd Y\bigr],
\]
the products being formed in the hatted variables, with finite fibres, because $\pi_m$ is a polynomial in $\zh$ at each $u$-order. The model residues are, term by term,
\[
\Res_{Y=0}\bigl[Y^m\re^{-W_V/\zh}\,\dd Y\bigr]=\sum_{i'\geq0}\frac{(-1/(4\zh))^{i'}}{i'!}\,\frac{(-V/\zh)^e}{e!}=s^{-(m+1)}\IC_m(v(u),z),\qquad e=m+4i'+1:
\]
the substitution $p=sY$, with $v=uV$ and $z=s^4\zh$, turns $\Wv(p)/z$ into $W_V(Y)/\zh$, and expanding $V^e$ in $u$ is a relabelling of the kind treated in \cref{lem:con-relab}. The exponent map $\iota$ of \cref{lem:con-relab} is additive, so the image of a product with finite fibres is the product of the images. Hence $\cV=s\sum_m\pi_m\,s^{-(m+1)}\IC_m=\Sigma$, the formal sum of \cref{lem:con-key}, coefficient by coefficient. By \cref{lem:con-id}, $\pi_{\nu_b}=\varepsilon_b\varpi_0G_b=\varepsilon_b\varpi_0s^{\nu_b}P_{0b}$, and $\YC_b=\varepsilon_b\IC_{\nu_b}$ (\cref{def:gau-van}); the signs cancel, $\varepsilon_b^2=1$, and so do the powers of $s$, $s\cdot s^{-(\nu_b+1)}\cdot s^{\nu_b}=1$; so $\Sigma=\varpi_0\sum_bP_{0b}\YC_b=\varpi_0(P\YC)_0$. No power of $\tau$ survives: the residue side carries none, and $G_b$ is defined without it. Since $\cV\in z^{-1}\bbQ[z^{-1}]\fps{u}$ (\cref{def:con-PhiSP}), the identity holds in $\bbQ\lser{z}\fps{u}$ (also by \cref{lem:con-weight}).
\end{proof}

\begin{thm}[Theorem P(b)]
\label{thm:con-Pb}
In every conifold frame of \cref{def:set-frames}, $\PhiSP=\Phican$ in $\bbQ\fps{u}\fps{\zh}$.
\end{thm}

The proof uses: Theorem~U (\cref{thm:gau-U}) for the pair of systems of \cref{prop:gau-quintic,prop:gau-equalprod}, whose hypotheses hold in the conifold frames; the gauge identity \cref{thm:gau-G}, with \cref{lem:gau-uniq,lem:gau-R1,lem:gau-vanmod,lem:gau-vanq}; \cref{cor:gau-cycsys}; the uniqueness half of \cref{prop:set-modrows}; \cref{lem:gau-poly}; \cref{thm:con-Pa}; \cref{lem:con-form,lem:con-cfu,lem:con-red,lem:con-S2,lem:con-id,lem:con-relab,lem:con-key,prop:con-res}; from \cref{sec:fsp}, \cref{lem:fsp-ibp,lem:fsp-cov,lem:fsp-cfu,lem:fsp-dq,lem:fsp-res,lem:fsp-stirfe,lem:fsp-transl,lem:fsp-homog}; the normalisations $\varpi_0(0)=1$ and $\tc'(0)=1$ of the conifold frames; and $\mathrm{ph}^{\mathrm{can}}_0=1$ (\cref{lem:gau-diag}, through \cref{thm:gau-R0}~(ii)). It uses neither (QRR), nor the finite route of \cref{sec:con-finite}, nor the degree bound \cref{prop:con-deg}, nor Theorem Gamma.

\begin{proof}
Steps 1--4 are \cref{lem:con-cfu} (the uniformising map), \cref{lem:con-red} (the reduction modulo exact forms), the invariance of both functionals under the uniformising map and their vanishing on exact forms (\cref{lem:fsp-cfu,lem:fsp-ibp,lem:fsp-res}), and the identification of \cref{lem:con-S2,lem:con-id}.

\emph{Step 5 (regularity).} By \cref{thm:con-Pa}~(iii), $\PhiSP$ solves HPF; by \cref{thm:gau-R0}~(ii), which rests on \cref{cor:gau-cycsys}, so does $\Phican$; both have constant term $1$ and coefficients in $\bbQ\fps{u}$. By \cref{lem:gau-poly}, $\PhiSP=\Phican\,\mathsf s(z)$ for a series $\mathsf s\in1+z^5\bbQ\fps{z^5}$ (the letter $\mathsf s$ is local to this proof; $z^5=u^4\zh^5$). Put $P^{\mathrm{SP}}:=\mathsf s(z)\Pc$. By \cref{thm:gau-U}~(b) and \cref{prop:gau-graded} it is a graded regular gauge, and $\Phi_{P^{\mathrm{SP}}}=\mathsf s(z)\Phican=\PhiSP$. The gauge identity for the canonical gauge, $\Pc\YC(v(u),z)=\YCQ$ (\cref{thm:gau-G}), multiplied by $\mathsf s(z)$, gives $(P^{\mathrm{SP}}\YC)_0=\mathsf s(z)(\YCQ)_0$.

\emph{Step 6 (the residue).} By \cref{prop:con-res} for $P=P^{\mathrm{SP}}$, $\cV=\varpi_0\,\mathsf s(z)\,(\YCQ)_0$ in $\bbQ\lser{z}\fps{u}$.

\emph{Step 7 (the order $u^1$).} On the left only $k=1$ contributes to $u^1$ in \cref{def:con-PhiSP}, because $\Xi=u+O(u^2)$, and $[x^0](\re^{-\cG/z}\fa)=1$, because $\cG=O(x^4)$ and $\fa(0,z)=1$; so $[u^1]\cV=-1/z$. On the right, $[u^0](\YCQ)_0=0$, since $y^{(0)}=-e_1$ and the terms $y^{(n)}$, $n\geq1$, vanish at $u=0$; and $[u^1](\YCQ)_0=-1/z$, since by \cref{lem:gau-vanq} only $y^{(1)}_0=-\tc$ contributes at $u^1$ and $\tc=u+O(u^2)$. With $\varpi_0(0)=1$ the right side has $u^1$-coefficient $-\mathsf s(z)/z$. Hence $\mathsf s=1$, and $\PhiSP=\Phican$.
\end{proof}

\begin{cor}[Frame independence]
\label{cor:con-frame}
$\Phican$ does not depend on the conifold frame.
\end{cor}

\begin{proof}
$\PhiSP$ is defined without reference to a frame (\cref{def:con-PhiSP}), and $\Phican=\PhiSP$ in every conifold frame by \cref{thm:con-Pb}.
\end{proof}

\begin{rmk}[The number $7/10$ from the residue]
\label{rmk:con-710}
At $u^2$, \cref{prop:con-res} for $P=\Pc$ reads $[u^2]\cV=[u^2]\bigl(\varpi_0(\YCQ)_0\bigr)$. On the left, the term $k=1$ of \cref{def:con-PhiSP} gives $-\frac1{2z}$, from $\Xi=u+\frac12u^2+\dots$, and the term $k=2$ gives $\frac1{2z^2}[x^1]\fa=-\frac1{5z}$, because the Stirling term $k=1$ of $\log\fa$ is $\bk_1(xz)(\frac15-5)=-\frac25xz$ with $\bk_1=\frac1{12}$. On the right only $z^{-1}$ occurs at $u^2$, with coefficient $-[u^2](\varpi_0\tc)=-[u^2]\varpi_1=-\gamma$ (\cref{lem:gau-vanq,lem:set-frobbasis}). So $-(\frac12+\frac15)/z=-\frac7{10}/z$: the Bernoulli number $B_2$ produces the coefficient $\gamma=\frac7{10}$ of the vanishing period. The remark is not used.
\end{rmk}

% Section 8 (part 8 of 8): the glue; the finite route.

\subsection{The glue}
\label{sec:con-glue}

For a colour-$0$ WKB function with polynomial coefficients write $\Phi(z,1):=\sum_n\mathrm{ph}_n(1)z^n$, its value at $u=1$, where $\zh=z$.

\begin{thm}[The glue]
\label{thm:con-glue}
In $\bbQ(\zeta)\fps{z}$,
\begin{equation}
\label{eq:con-glue}
h(z)\,\Dn_0(z)=\Phican(z,1);
\end{equation}
hence $\Dn_0\in1+z\bbQ\fps{z}$, and, by \cref{thm:con-Pb}, $\Dn_0(z)=\PhiSP(z,1)/h(z)$.
\end{thm}

\begin{proof}
We recall how \cref{thm:gau-R0}~(ii)--(iv) places both columns in one ring. (a) The canonical flat amplitude $A^{\mathrm{can}}=\sum_nA^{\mathrm{can}}_nz^n$ has $A^{\mathrm{can}}_n=L^{4n+1}\mathrm{ph}^{\mathrm{can}}_n(L^{-5})\in\bbQ[L,L^{-1}]$, the $\mathrm{ph}^{\mathrm{can}}_n$ being polynomials. (b) In $\bbQ(\zeta)\fps{q}$, with $L=(1-5^5q)^{-1/5}$, $A^{\mathrm{can}}$ and the actual amplitude $A^{\mathrm{act}}$ are colour-$0$ WKB columns of $\PFfive$ with the same exponential $\re^{u_0/z}$ and the same leading term $L$; by formal WKB uniqueness (\cref{lem:gau-wkb}, with $\Bbbk=\bbQ(\zeta)$) they differ by a constant series in $1+z\bbQ(\zeta)\fps{z}$. (c) Hence every $A^{\mathrm{act}}_n$ lies in $\bbQ(\zeta)[L,L^{-1}]$. (d) The ring $\bbQ(\zeta)[L,L^{-1}]\fps{z}$ embeds in $\bbQ(\zeta)\lser{s}\fps{\zh}$, where the comparison of \cref{thm:gau-R0} takes place, with the same exponential stripped on both sides; there $A^{\mathrm{act}}\,\Dn_0=A^{\mathrm{can}}$, the constant series of (b) is $\Dn_0\in1+z\bbQ(\zeta)\fps{z}$, and it does not depend on $u$ (the actual function has constant term $1$, so it is not a zero divisor). (e) Evaluation at $L=1$ is a ring homomorphism $\bbQ(\zeta)[L,L^{-1}]\to\bbQ(\zeta)$; applied coefficientwise it gives $A^{\mathrm{act}}(z)|_{L=1}\,\Dn_0(z)=A^{\mathrm{can}}(z)|_{L=1}$ in $\bbQ(\zeta)\fps{z}$. Since $A^{\mathrm{can}}_n|_{L=1}=\mathrm{ph}^{\mathrm{can}}_n(1)$, the right side is $\Phican(z,1)$, a series with rational coefficients. The point $L=1$ is $u=1$, $q=0$, where $\zh=z$; as $(1-5^5q)^{-1/5}=1+O(q)$, the value at $L=1$ of the Laurent polynomial $A^{\mathrm{act}}_n$ is the coefficient of $q^0$ of its image in $\bbQ(\zeta)\fps q$, and $A^{\mathrm{act}}_n(q=0)=h_n$ for every $n$, by (QRR) in the form of \cref{thm:set-qrr}, \eqref{eq:set-qrr}; so the left side is $h(z)\Dn_0(z)$. As $h\in1+z\bbQ\fps{z}$ is invertible in $\bbQ\fps{z}$, $\Dn_0=\Phican(z,1)/h(z)$ has rational coefficients, and \cref{thm:con-Pb} gives the second form. In particular the actual amplitude has rational coefficients, $A^{\mathrm{act}}_n=\sum_{m\leq n}A^{\mathrm{can}}_m\coef{z^{n-m}}\Dn_0^{-1}\in\bbQ[L,L^{-1}]$. This is the only place in this section where an input enters: (QRR), together with the actual rows and their WKB form, which rest on (M) and (G0) (\cref{thm:set-mirror,thm:set-G0}) through \cref{cor:set-cyclicPF5} and \cref{thm:gau-R0}. Neither the finite route of \cref{sec:con-finite} nor the degree bound \cref{prop:con-deg} is used.
\end{proof}

\begin{cor}
\label{cor:con-A2}
$\Dn=1$ if and only if $\PhiSP(z,1)=h(z)$. By \cref{thm:con-Pa}~(iv), $\PhiSP(z,1)=c(z;1)$, the value at $L=1$ of the Laurent polynomials $c_n$; so the condition reads $c(z;1)=h(z)$ (\cref{cor:gam-equiv}).
\end{cor}

\begin{proof}
By \cref{cor:gau-colours}, $\Dn_i(z)=\Dn_0(\zeta^{-i}z)$, so $\Dn=1$ if and only if $\Dn_0=1$; by \eqref{eq:con-glue} and \cref{thm:con-Pb}, $\Dn_0=1$ if and only if $\PhiSP(z,1)=h(z)$. By \cref{thm:con-Pa}~(iv), $c_n(1)=\mathrm{ph}^{\mathrm{SP}}_n(1)$.
\end{proof}

The condition of \cref{cor:con-A2} is a connection statement: $\PhiSP$ is built at the conifold point $u=0$, as a series in $u$ whose coefficients turn out to be polynomials, and it is compared at $u=1$ with the value prescribed there. The gauge, the periods and the model have dropped out of it.

\begin{rmk}[Why an exact input is needed]
\label{rmk:con-why}
Any two colour-$0$ WKB columns of $\PFfive$ with the same exponential differ by a constant series (\cref{lem:gau-wkb}). The properties used so far, namely $\PFfive$, symplecticity (which could at most give $\Dn_0(z)\Dn_0(-z)=1$, and only through the symplecticity of the expanded rows, \cref{rmk:con-genus}), the $\bbZ/5$-grading and the polynomial structure (\cref{lem:gau-poly,prop:con-deg}), are all invariant under $\Dn_0\mapsto\Dn_0\exp(a(z))$ for every odd series $a\in z^5\bbQ\fps{z^{10}}$: multiplication by such a factor commutes with HPF, keeps the coefficients polynomial ($z^5=u^4\zh^5$) and preserves $\Dn_0(z)\Dn_0(-z)=1$. So no argument using only these properties decides $\Dn_0$. Theorem~P removes the conifold end by an exact identity, the residue at the pole of the $1$-form (\cref{prop:con-res}), and Theorem Gamma supplies the exact input at $q=0$ (\cref{sec:gam}); both arguments are formal, but their input is the Mellin--Barnes integrand. Among the conclusions of the natural factorisation \cref{cor:mai-natfact}, the symplecticity of $G$ is the only one that a change of the rows by a scalar factor $\cn(z)^{-1}$, $\cn\in1+z^5\bbQ\fps{z^5}$, can violate: the factor multiplies $G$ by $\cn(z)^{-1}$, which is $1$ at $u=0$ and at $\zh=0$ and has polynomial coefficients ($z^5=u^4\zh^5$), and $G\cn^{-1}$ is symplectic if and only if $\cn(z)\cn(-z)=1$; so it is blind to the factors $\exp(a(z))$ with $a$ odd. The computations of \cref{sec:evi-normal} illustrate this; they are not used.
\end{rmk}

\subsection{The finite route}
\label{sec:con-finite}

This subsection reads the normalisation at the conifold point order by order. It explains the computations of \cref{sec:evi-normal} and is not used by the proof of the main theorem. Let $\Phi$ be a colour-$0$ WKB function solving HPF with coefficients in $\bbQ\fps{u}$; by \cref{lem:gau-poly} and \cref{thm:gau-R0}~(ii), $\Phi=\Phican\,\cn(z)$ for a unique $\cn\in1+z^5\bbQ\fps{z^5}$. Its hatted rows $\hat\Qr(\Phi)$ are built as follows: row $0$ of colour $0$ is $\Phi/\varpi_0$; the rows $a=1,\dots,4$ follow by the forward relations of the cyclic system (\cref{cor:gau-cycsys}); colour $i$ follows from $\Phi_i(\zh,u)=\Phi(\zeta^{-i}\zh,u)$ and the covariance $[\zh^n]\hat\Qr^{(i)}_a=\zeta^{i(a+1-n)}[\zh^n]\hat\Qr^{(0)}_a$ (\cref{lem:gau-colrows}). Put $\Ghat(\Phi):=\hat\Qr(\Phi)\,\Rhat^{\mathrm m}(\tau\zh)^{-1}$.

\begin{lem}[Reading the normalisation]
\label{lem:con-N}
$\Ghat(\Phi)$ is the hatted form $\Lambda_sP(u,s^4\zh)\Lambda_\varsigma^{-1}$, $\Lambda_s=\diag(s^{1-\nu_a})$, $\Lambda_\varsigma=\diag(\varsigma^{1-\nu_b})$, of the graded regular gauge $P=\Pc\,\cn(z)$, and $P_{11}(0,z)=P_{22}(0,z)=\cn(z)$. Moreover $[u^j\zh^{5r}]\Ghat_{11}=0$ for $j<4r$, and
\[
\mathsf K_r:=[u^{4r}\zh^{5r}]\Ghat_{11}=[z^{5r}]\cn=[u^{4r}\zh^{5r}]\Ghat_{22}\qquad(r\geq1).
\]
As a function of the free coefficients of $\Phi$, $\mathsf K_r-\fc_r$ depends only on $\fc_1,\dots,\fc_{r-1}$. (The index $r$ is a local letter.)
\end{lem}

\begin{proof}
For $\Phi=\Phican$ the rows built in this way are the hatted rows of $\Pc\Qr^{\mathrm m}(v(u))$: this matrix solves the quintic system (\cref{thm:gau-R0}~(i)), so its rows follow from row $0$ by the forward relations, and it is colour covariant (\cref{lem:gau-colmod,prop:gau-graded}). The forward relations and the colour covariance commute with multiplication by $\cn(z)$, because $\cn(\zeta z)=\cn(z)$; so $\hat\Qr(\Phi)=\cn(z)\,\hat\Qr(\Phican)$ and $\Ghat(\Phi)$ is the hatted form of $\Pc\cn(z)$, which is a graded regular gauge (\cref{thm:gau-U}~(b), \cref{prop:gau-graded}), with $P_{11}(0,z)=\cn(z)$ since $\Pc$ is canonical. By \cref{lem:gau-R1}, $P_{22}(0,z)=P_{11}(0,z)$. For $b\in\{1,2\}$ the hatted diagonal entry is $\Ghat_{bb}=\tau^{(1-\nu_b)/4}P_{bb}(u,s^4\zh)$, and the graded entry $P_{bb}$ contains only $z^{5r}=u^{4r}\zh^{5r}$, with coefficient $p_{5r,b}(u)$; so $[\zh^{5r}]\Ghat_{bb}=\tau^{(1-\nu_b)/4}u^{4r}p_{5r,b}(u)$, whose $u^{4r}$-coefficient is $p_{5r,b}(0)=[z^{5r}]P_{bb}(0,z)=[z^{5r}]\cn$. Finally $\fc_r=[u^{4r}\zh^{5r}](\Phican\cn)=\sum_{m=0}^r[z^{5m}]\cn\cdot\fc^{\mathrm{can}}_{r-m}$, where $\fc^{\mathrm{can}}_k:=[u^{4k}]\mathrm{ph}^{\mathrm{can}}_{5k}$ and $\fc^{\mathrm{can}}_0:=1$, by extracting the coefficient of $u^{4r}\zh^{5r}$ from the product with $\cn(z)=\sum_m[z^{5m}]\cn\,u^{4m}\zh^{5m}$; so $\mathsf K_r=\fc_r-\fc^{\mathrm{can}}_r-\sum_{0<m<r}[z^{5m}]\cn\cdot\fc^{\mathrm{can}}_{r-m}$, and the $[z^{5m}]\cn$, $m<r$, are determined by $\fc_1,\dots,\fc_{r-1}$ in the same way.
\end{proof}

Let $\Phi^{\mathrm{pol}}$ be the solution of HPF of \cref{prop:con-deg} whose free coefficients are fixed by $\mathrm{ph}^{\mathrm{pol}}_{5m}(1)=h_{5m}$, $m\geq1$; this is possible because $\fc_m$ enters $\mathrm{ph}_{5m}(1)$ with coefficient $1$ and does not enter $\mathrm{ph}_n$ for $n<5m$.

\begin{thm}[The connection form]
\label{thm:con-conn}
With $\Phi^{\mathrm{pol}}$ as above:
\begin{enumerate}
\item $\Dn_0\in\bbQ\fps{z^5}$ if and only if $\mathrm{ph}^{\mathrm{pol}}_n(1)=h_n$ for every $n$ with $5\nmid n$.
\item If $\Dn_0\in\bbQ\fps{z^5}$, then the actual colour-$0$ function is $\Phi^{\mathrm{pol}}$, $\Dn_0=\cn^{-1}$ for the series $\cn$ of \cref{lem:con-N} applied to $\Phi^{\mathrm{pol}}$, and $\Dn_0=1$ if and only if $\mathsf K_r=0$ for every $r\geq1$; the conditions $\mathsf K_r=0$, $r\geq1$, characterise the free coefficients of $\Phican$.
\item $\Dn_0=1$ if and only if $\mathrm{ph}^{\mathrm{can}}_n(1)=h_n$ for every $n$.
\end{enumerate}
\end{thm}

\begin{proof}
(The letters $\mathsf p$, $\mathsf d$ are local to this proof.) By \cref{lem:gau-poly}, $\Phi^{\mathrm{pol}}=\Phican\,\mathsf p(z)$ with $\mathsf p\in1+z^5\bbQ\fps{z^5}$, and by \cref{thm:gau-R0}~(iv) the actual function is $\Phican/\Dn_0=\Phi^{\mathrm{pol}}/\mathsf d$ with $\mathsf d:=\mathsf p\Dn_0$; by \eqref{eq:con-glue}, $\Phi^{\mathrm{pol}}(z,1)=h(z)\,\mathsf d(z)$. (i) $\Dn_0\in\bbQ\fps{z^5}$ if and only if $\mathsf d\in\bbQ\fps{z^5}$. If $\mathrm{ph}^{\mathrm{pol}}_n(1)=h_n$ for every $n$ with $5\nmid n$, then $\Phi^{\mathrm{pol}}(z,1)=h(z)$, the multiples of $5$ agreeing by construction, so $\mathsf d=1$. Otherwise let $N$ be the least $n$ with $\mathrm{ph}^{\mathrm{pol}}_n(1)\neq h_n$; then $5\nmid N$, and $\mathsf d=1+(\mathrm{ph}^{\mathrm{pol}}_N(1)-h_N)z^N+O(z^{N+1})$ is not in $\bbQ\fps{z^5}$. (ii) By (i), $\mathsf d=1$, so the actual function is $\Phi^{\mathrm{pol}}$ and $\mathsf p=\Dn_0^{-1}$; \cref{lem:con-N} applied to $\Phi^{\mathrm{pol}}$ has $\cn=\mathsf p$. Hence $\Dn_0=1$ if and only if $\cn=1$, if and only if $\mathsf K_r=[z^{5r}]\cn=0$ for every $r$, and by the last clause of \cref{lem:con-N} these conditions fix the free coefficients recursively, at the values of $\Phican$. (iii) is \eqref{eq:con-glue}.
\end{proof}

By \cref{cor:con-frame}, $\mathsf K_r=[z^{5r}](\Phi/\Phican)$ and $\Dn_0=\Phican(z,1)/h(z)$ do not depend on the conifold frame; computations of $\mathsf K_r$ in several frames (\cref{sec:evi-normal}) are therefore no independent evidence for $\Dn_0=1$.


